\scanpage{5}
Plus précisément, le foncteur composé
\diagram{\begin{tikzcd}[ampersand replacement=\&]
D_{\mathrm{parf}}(A)\arrow[r,"\Phi_L"] \& D_{\mathrm{parf}}(B)\arrow[r,"{\Phi_{{}^{s}L}}"] \& D_{\mathrm{parf}}(A)
\end{tikzcd}}
est \inserted{can.} isomorphe au foncteur \deleted{identique, \unsure{par suite}} \inserted{$\mathbb{L} s_A^*\deleted{\unsure{Rs}}\simeq R s_{A*}$ suivi de $\tau_A[-n]$}, \deleted{\unsure{ce qui|de sorte que|respectivement}} où $\tau_A\simeq\unsure{\bigwedge^{n}|\bar{\Lambda}|\hat{\Lambda}}t_A\simeq\bigwedge^{n}t_B$. Donc $\Phi_L$ et $\mathbb{L} s_A^*\Phi_{{}^{s}L}\inserted{=\mathbb{R} s_{A*}\Phi_{{}^{s}L}}$ sont des équivalences des catégories triangulées de $D_{\mathrm{parf}}(A)$ et $D_{\mathrm{parf}}(B)$.
% slanted
\margin{inutile de se borner aux complexes parfaits, on peut les prendre p.ex. dans $D^b_{\mathrm{qcoh}}\deleted{\unsure{A}}$}
\margin{\deleted{\unsure{des catégories|identique|identité}}}

\uline{Corollaire} \unsure{3|2} $\Phi_L\circ\omega_A[n]$ : \deleted{transforme} \inserted{$F\longmapsto\Phi_L(\unsure{F|\omega_A}\otimes\omega_A[n])\simeq\Phi_L(F)\otimes\omega_A[n]$ transforme} $\overset{\mathbb L}{\otimes}$ en $\overset{\mathbb L}{*}$ : \deleted{(et \unsure{inversement} $\overset{\mathbb L}{*}$ en $\overset{\mathbb L}{\otimes}$)}. D'où $=$ \deleted{$\unsure{(-1)^n\varphi_L}$} \deleted{$\unsure{x\mapsto\varphi_L(x}$} \inserted{$x\longmapsto(-1)^n\mathrm{cl}\,\omega_A\,\varphi_L(x)$} transforme produit ordinaire en \deleted{\unsure{cup}} produit de convolution.
\[
\begin{aligned}
\Phi_L(F\unsure{\overset{\mathbb L}{\otimes}|\otimes}G)\otimes\omega_A[n]&\simeq[\Phi_L(F)\otimes\omega_A[n]]*[\Phi_L(G)\otimes\omega_A[n])\\
\text{i.e.}\quad \Phi_L(F\overset{\mathbb L}{\otimes} G)&\simeq[\Phi_L(F)*\Phi_L(G)]\unsure{\otimes|\mathbin{\cdot}}\omega_A[n]
\end{aligned}
\]

Prouvons la prop. 2. Il suffit de définir un isom $\Phi_{{}^{s}L}\circ\Phi_L\simeq \mathrm{id}_{D_{\mathrm{parf}}(A)}$. Or le composé des foncteurs est de la forme $F\longmapsto \mathbb{R} \mathrm{pr}_{2*}(\mathrm{pr}_1^*(F)\overset{\mathbb L}{\otimes} M)$, où $M\in D_{\mathrm{parf}}(A)$ est donné par
\margin{$A\times B\times A$}
\[
M=\unsure{\mathbb{R}|\mathbb L R}\mathrm{pr}_{13*}\bigl(\mathrm{pr}_{12}^*(L)\,\mathrm{pr}_{23}^*({}^{s}L)\bigr)
\]
Il faut donc calculer ce $\mathbb{R} \mathrm{pr}_{13*}$, \deleted{et} Pour ceci, posons $C=A\times A$, et considérons l'isom $A\times B\times A\simeq\deleted{\unsure{C\times B\times}}C\times B$. Le module inversible $N=\mathrm{pr}_{12}^*(L)\,\mathrm{pr}_{23}^*({}^{s}L)$ sur $C\times B$ est \deleted{\unsure{muni de}} birigidifié de façon canonique, et correspond à un homomorphisme $C\xrightarrow{\psi}B^*$, qui n'est autre que l'hom. de $A\times A$ dans $B^*$ dont les \inserted{deux} composantes sont égales à $\varphi\colon A\unsure{\mathrel{-}|\to}B$ (l'hom de $A$ dans $B^*$ défini par $L$).
% archive sheet number: 4
